\documentclass[10pt,a4paper]{amsart}
\usepackage[margin=19mm]{geometry}
\usepackage{amsmath,amssymb,mathtools}
\usepackage[scr=rsfs,cal=euler]{mathalfa}
\usepackage{xfrac}
\usepackage[unicode,hidelinks,bookmarks=false]{hyperref}
\newcommand{\ie}{i.e.\ }
\newcommand{\cf}{cf.\ }
\newcommand{\resp}{respectively\ }
\newcommand{\Subsection}{\subsection}
\providecommand{\nobreakdash}{\nobreak\hbox{-}}
\setlength{\emergencystretch}{2em}
\multlinegap0pt
\usepackage{float}
\usepackage{amssymb}
\usepackage{xcolor}
\usepackage[shortlabels]{enumitem}
\usepackage{tikz}
\usepackage{cancel}
\newtheorem{lemma}{Lemma}
\usepackage{cleveref}
\crefname{lemma}{Lemma}{Lemmas}
\newcommand{\IT}{\mathrm{LPF}} 
\newcommand{\Out}{\mathcal{O}}
\newcommand{\arms}{\mathrm{Arms}}
\newcommand{\legs}{\mathrm{Legs}}
\title{Lehmer Parking Functions and Their Outcomes}
\author{Melissa Beerbower, Jennifer Elder, Pamela E. Harris, Ilana Lavene, Lucy Martinez, Adam Martinson, Molly Oldham}
\date{Source: arXiv:2603.20535v1 (CC BY 4.0)}
\begin{document}
\maketitle

\noindent\emph{Source and licence.} Lehmer Parking Functions and Their Outcomes. Version arXiv:2603.20535v1 (CC BY 4.0), distributed by arXiv under the Creative Commons Attribution 4.0 International licence, http://creativecommons.org/licenses/by/4.0/, as recorded in the arXiv licence metadata for this exact version and shown on the abstract page https://arxiv.org/abs/2603.20535v1. Full licence notice: \texttt{licenses/arxiv-2603.20535-CC-BY-4.0.txt}.\par\smallskip

\noindent\textit{Editorial scope.} Counted unit: the article's lemma surjective (source0.tex lines 647--650). The article's complete original proof of it is delivered with it (source0.tex lines 651--669, one \texttt{proof} environment of 19 lines). No mathematics is written, repaired or re-derived: the statement and the proof are the article's own lines, in the article's own notation, and no retained line is substituted or re-flowed.  Retained uncounted as the source-local context the proof needs: lines 520--522; lines 578--580; lines 615--620; lines 671--676.  One declared adaptation: exactly 2 ancillary strings inside the retained spans are omitted (image-inclusion commands and, where the source repeats a label, the repeated label definitions) and no other line differs from the source; the figure environments, with their placement, captions and labels, are kept, and the package records the omitted strings in fidelity receipts bound to the affected bodies.  No retained line contains a citation, so the wrapper carries no bibliography and no bibliography entry.  The retained lines contain the article's own diagram code, which the wrapper typesets with the standard tikz packages; no image file is delivered and the package declares no asset dependency.  Editorial formatting only, in addition: an AMS \texttt{amsart} preamble that restates the article's own declarations and macros used by the retained lines, namely the environments \texttt{lemma} on separate counters, together with the standard packages those lines need.  The retained environments print in this document as Lemma 1 in order of appearance, whereas the article prints the same items with the numbering of its own full document.  The complete native source of the article as distributed in the arXiv e-print for this exact version is bundled unchanged as \texttt{sources/196856/source0.tex}.
\begin{lemma}\label{lem:d_i = leg intersections}
    Let $\pi\in \Out (\IT_n)$ and consider the spaced parenthesization $(\arms_\pi,\legs_\pi)$. Then for all $i\in [n]$, $d_i$ is the number of peaks of the arm-leg diagram of $\pi$ lying in the box $[i,n]\times[n-i+1,n]$. Consequently, $d_i$ is nonnegative for all $i\in[n]$.
\end{lemma}

\begin{lemma}\label{lem:balanced -> non-intersecting partial}
    If $(F,L)$ is a balanced spaced parenthesization of $n$ spaces, then there exists a unique non-intersecting partial arm-leg diagram $\tau$ of the grid $[n]\times[n]$ such that $(\arms_\tau,\legs_\tau)=(F,L)$. If $(F,L)$ has $k$ pairs of parentheses, then $\tau$ has $k$ entries. 
\end{lemma}

\begin{figure}[H]
    \centering
    
    \caption{In \Cref{lem:balanced -> non-intersecting partial} (left), we place peaks in rows $i$ satisfying $n-i+1\in F$. In \Cref{lem:Phi is surjective} (right), we place the remaining entries below the antidiagonal. Here, there are four possible outcomes of Lehmer parking functions with the given parenthesization.}
    \label{fig:algorithm for below antidiagonal}
\end{figure}

\begin{figure}[H]
    \centering
    
    \caption{We assume that row $n-i+1$ is the uppermost empty row. There are $i-1$ entries to the left of the antidiagonal entry $(i,n-i+1)$, and $(i-1)-|\text{box }A|$ are available for the next entry (\cref{eq:2}). Likewise, there are $i-1$ entries in the top $i-1$ columns, and $(i-1)-|\text{box }A|$ of them lie in box $B$ (\cref{eq:3}). Finally, row $n-i+1$ is empty, so the same number of entries lie in box $C$ (\cref{eq:4}), the box given in \Cref{lem:d_i = leg intersections}.}
    \label{fig:proof_3.6}
\end{figure}

\begin{lemma}\label{lem:Phi is surjective}
    If $(F,L)$ is a balanced spaced parenthesization, then there exists at least one 
    $\pi \in \Out(\IT_n)$ such that $(F,L)=(\arms_\pi,\legs_\pi)$.
\end{lemma}

\begin{proof}
    Given $(F,L)$, we start from the non-intersecting partial arm-leg diagram $\tau$ from \Cref{lem:balanced -> non-intersecting partial} satisfying $(\arms_\tau,\legs_\tau)=(F,L)$ (see \Cref{fig:algorithm for below antidiagonal}, left). Because $\arms_\tau=F$, we have that $\tau$ contains entries in each row $n-i+1$ such that $i\in F$.
    We now construct $\pi$ from $\tau$ as follows: we focus on the rows which do not contain an entry, and in each such row 
    beginning from top to bottom, iteratively place an entry $(j,n-i+1)$ below the antidiagonal, where $j\in[1,i-1]$ and $i\notin F$,
    (see \Cref{fig:algorithm for below antidiagonal}, right).
    We show that if $n-i+1$ is the uppermost empty row, 
    there exists at least one column numbered $j\in[1,i-1]$ which contains no entries of the diagram. This implies that we can place the point $(j,n-i+1)$ below the antidiagonal.
    Note that \Cref{lem:d_i = leg intersections} follows identically for $(\arms_\tau,\legs_\tau)$; that is, $d_i$ is given by the number of entries of $\tau$ in $[i,n]\times[n-i+1,n]$.
    We  illustrate the following equations in \Cref{fig:proof_3.6}, where the figure numbers correspond to the equation numbers below. In the caption we provide some additional details to explain the following count: 
    The number of empty columns $j\in[1,i-1]$ is given by 
    \begin{align}
        \label{eq:2}&(i-1)-\text{number of entries in }[1,i-1]\times[n-i+2,n]&\text{no lower entries below the antidiagonal}\\
        \label{eq:3}=\,&\text{number of entries in }[i,n]\times[n-i+2,n]&\text{the top $i-1$ rows have $i-1$ entries}\\
        \label{eq:4}=\,&\text{number of entries in }[i,n]\times[n-i+1,n]&\text{row $n-i+1$ is empty}\\
        =\,&d_i,&\text{by \Cref{lem:d_i = leg intersections} for partial diagrams}
    \end{align}
    which is positive because $(\arms_\tau,\legs_\tau)$ is balanced. 
    Placing entries $(j,n-i+1)$ in this manner for each remaining row yields an outcome of a Lehmer parking function, $\pi$, satisfying $(\arms_\pi,\legs_\pi)=(F,L)$.
\end{proof}

\end{document}
